\chapter{Lithotripsy}

\section*{Overview}
ESWL uses shock waves generated outside the body to break kidney or ureteral stones into tiny fragments (sand) that can then pass naturally in the urine.

\section*{Alternative Names}
Extracorporeal Shock Wave Lithotripsy, ESWL, Kidney stone blast

\section*{Principle}
Shock waves are generated outside the body and focused on a kidney or ureteral stone, traveling through tissue to produce compressive and tensile forces. These forces fragment the stone into small pieces that can pass in the urine.
\section*{History}
Developed in Germany by Dornier Medizintechnik in the late 1970s. The first human treatment occurred in 1980, transforming kidney stone management.

\section*{Why It Is Done}
Ordered for symptomatic kidney or ureteral stones, typically between 5 mm and 20 mm in size, that fail to pass spontaneously.

\section*{Is This a Primary or Secondary Test or Procedure}
Primary non-invasive treatment option for simple renal pelvic and upper ureteral stones.

\section*{How It Is Performed}
Performed under sedation or general anesthesia. You lie on a table. Fluoroscopy or ultrasound is used to locate the stone. A shock-wave generator delivers thousands of targeted shock waves through a water cushion.

\section*{Preparation}
Fast only as directed by the facility, especially if sedation or anesthesia is planned. Tell the treating team about every anticoagulant, antiplatelet medicine, and supplement. Do not stop these medicines unless the prescribing and procedural teams provide a drug-specific plan that balances bleeding and clotting risks.

\section*{Contraindications and Precautions}
Pregnancy, active urinary tract infection, uncorrected bleeding disorders, or severe obesity preventing stone targeting.

\section*{Risks}
Renal hematoma, hematuria (blood in urine, expected), renal colic (pain from passing fragments), or 'steinstrasse' (ureteral obstruction by stone fragments).

\section*{Aftercare and Recovery}
Drink plenty of water to flush out stone fragments. Strain your urine to collect fragments for analysis. Take prescribed alpha-blockers.

\section*{Outcome and Follow-up}
Success is determined by the fragmentation of the stone and complete passage of fragments, confirmed by follow-up X-ray or ultrasound.

\section*{Expected Results}
Complete fragmentation of the target stone; passage of stone gravel; no residual stones on post-op imaging.

\section*{Follow-up}
Follow-up X-ray (KUB) or ultrasound in 2 weeks; dietary counseling for stone prevention.

\section*{When to Seek Medical Care}
Follow the procedure team's specific instructions. Seek urgent medical assessment for severe or worsening pain, uncontrolled bleeding, fever or signs of infection, trouble breathing, chest pain, fainting, new weakness or confusion, inability to urinate, or any rapidly worsening symptom. The appropriate warning signs depend on the procedure and the patient's medical condition.

\section*{References}
\begin{enumerate}
  \item MedlinePlus. Lithotripsy. U.S. National Library of Medicine; 2025.
  \item Current Medical Diagnosis and Treatment. McGraw-Hill; 2025.
\end{enumerate}

\clearpage
